\documentclass[a4paper,11pt,openany]{ltjsbook}
\usepackage{luatexja}
\usepackage{luatexja-fontspec}
\usepackage{amsmath,amssymb,amsfonts,amsthm}
\usepackage{geometry}
\geometry{top=24mm,bottom=24mm,left=24mm,right=24mm,footskip=10mm}
\usepackage{tikz}
\usepackage{tcolorbox}
\tcbuselibrary{skins,breakable}
\usepackage{hyperref}
\usepackage{fancyhdr}

\definecolor{primaryblue}{RGB}{20, 60, 110}
\definecolor{softblue}{RGB}{238, 244, 252}
\definecolor{darkslate}{RGB}{35, 45, 60}

\hypersetup{
    colorlinks=true,
    linkcolor=primaryblue,
    citecolor=darkslate,
    pdftitle={古典で学び直すド・モルガン数学思考入門},
    pdfauthor={教育情報分析研究会}
}

\pagestyle{fancy}
\fancyhf{}
\fancyhead[LE,RO]{\thepage}
\fancyhead[RE]{\leftmark}
\fancyhead[LO]{\rightmark}
\setlength{\headheight}{14pt}
\renewcommand{\headrulewidth}{0.4pt}

\title{
    {\huge\bfseries\gtfamily 古典で学び直すド・モルガン数学思考入門}\\[14pt]
    {\large 底本：\textit{On the Study and Difficulties of Mathematics} (Project Gutenberg eBook \#39088)}
}
\author{\textbf{教育情報分析研究会} 編著}
\date{2026年10月}

\begin{document}
\maketitle
\frontmatter
\tableofcontents
\clearpage
\mainmatter

\chapter{Chapter I.: Introductory Remarks on the Nature and Objects of Mathematics.}

\#\#\# Chapter I: 数学の本質と目的について

この論文では、数学を学ぶために最適な方法や道程を示し、初学者が直面する困難な点について詳しく説明します。また、歴史的事実や形而上学的な主張とは異なる、数学の信頼性と証明の特性について深く掘り下げます。

まず、数学は他の科目とは異なる特異な性質を持っています。言語や歴史などの学問では、事実を確認するためには多くの手順が必要です。例えば、キーサーが殺害されたという事実は、多くの信頼できる人々の記録と証拠に基づいています。しかし、数学においては、その性質が全く異なります。

\#\#\#\# 数学における証明

数学では、ある命題が真であることを絶対的に証明します。これは歴史的事実とは異なる点です。例えば、「キーサーはブロツとその友人によって暗殺された」という事実は、その出来事が実際に起こったという証拠があるからこそ信頼できます。しかし、数学における証明は、ある命題が無理で不可能であることを示すことです。

\#\#\#\# 几何学の基本原理

幾何学では、以下の3つの基本原理があります：

1. **量を部分に分割した場合、全体はどの部分よりも大きい**（I）
2. **二つの直線は空間を囲むことができない**（II）
3. **一点を通じてはただ一つの直線が、他の直線と交わらないように描くことができる。そのような直線は「平行」である**（III）

これらの原理に基づいて、幾何学全体が構築されます。各命題の証明は、それらの原理と矛盾しないことを示すことで成り立ちます。例えば、ユークリッド『原論』第1巻、命題4では、二つの三角形がそれぞれ二辺とその挟角が等しい場合、これらはすべてにおいて等しいことを証明しています。これは、もし等しくないならば、二つの直線が空間を囲むことになるという矛盾を示します。

\#\#\#\# 感覚と論理

数学の最初の教師は私たちの感覚です。数や点、直線、面などの概念は、単語で表現するしか方法はありません。しかし、これらの概念は非常に明確で、10歳児でも理解できるほどです。さらに、これらの概念を論理的に考察するための手段も提供されます。

例えば、「全体はその一部より大きい」という命題は、基本的な事実として提示されます。これは「二加二が四」のように、基本的な真理から有用な結論を得ることができます。数学初心者は、このような論理的真理から有用な結論を得ることは不可能だと感じるかもしれませんが、それは誤解です。

\#\#\#\# 進歩の難しさ

幾何学は最も単純な開始点を提供しますが、進行するにつれて困難が急速に増大します。初めから見れば多少パラドックスに思えるかもしれませんが、より鋭い学生ほど進歩の妨げになります。しかし、これらの基本原理を理解し、自身で体現することが非常に重要です。

\#\#\#\# 結論

数学は他の学問とは異なる性質を持ち、その証明方法も独特です。歴史的事実や形而上学的な主張とは異なり、数学の信頼性は絶対的な真理に基づいています。初学者にとって、これらの基本原理を理解し、自身で体現することが非常に重要です。

このように、数学の学習は、単に知識を吸収するだけでなく、論理的思考と証明の能力を鍛えるプロセスでもあります。その過程を通じて、私たちが持つ感覚と論理的な考察方法を深く理解し、より深い洞察を得ることができます。

この学びは、私たちの心に新しい視点をもたらし、数学の美しさと奥深さを体験させてくれます。ぜひ、この旅を一緒に楽しみましょう。

\chapter{Chapter II.: On Arithmetical Notation.}

\#\#\# Chapter II.: On Arithmetical Notation

初等的な算術の概念は、他の科学も同様に、初期の観察から導き出されるものです。これらの概念がどのように心に生まれるかについては探求する必要はありませんし、数や量という言葉の意味を定義することも不可能です。これらは非常に単純な概念であり、つまりそれらを表すアイデアが私たちの心の中での最も基本的なアイデアだからこそ、よりシンプルな他の概念で説明することは不可能であると言えます。これが用語を定義するという意味であり、ここでは一般の定義について少し述べることで、すべての科学に適用できることが言えるでしょう。

定義は、他のより容易に理解できる用語を使用してある用語を解説し、その意味を固定することで概念を明確にするものです。これにより、その用語が含意するものは何か、また含意しないものは何かがはっきりと見えてきます。定義自体が証明すべき事実や他の仮定を内包していないか十分に注意しなければなりません。例えば、「正方形とは四角形であり、すべての辺が等しく、すべての角度が直角である」という定義では、正方形の真実性は述べられていますが、必要以上の情報を含んでいます。なぜなら、四角形のすべての辺が等しい場合でも、一つだけの角度が直角であれば他のすべての角度も直角であることが証明できるからです。したがって、上記の定義では、証明すべき事実を暗に含んでいることになります。

また、上記の定義は一方で冗長であり、他方で欠落していると言えます。なぜなら、図形の辺が直線か曲線かについて述べられていないからです。正しくは、「正方形とは四角形であり、すべての辺が等しく、一つの角度が直角である直線的な四角形」であるべきです。

数学の科学は、その証明の優越性において、用語がどれだけ正確に定義されているかに大きく依存しています。初学者が良い定義とは何かを明確に理解することほど重要なものはありません。数や量のような定義できない用語があることは認めますが、これらの用語に対する定義は学生に困難をもたらすばかりでなく、既に幾何学では点や直線などの用語の説明が試みられています。点は「部分を持たないものであり、大きさを持たないもの」と定義されます；直線は「極端な点間に均等に配置されたもの」です。

ここで誰でも自問してみてください：もし幾何学を始める前に「部分を持たないものであり、大きさを持たないもの」と「極端な点間に均等に配置された線」という言葉について話されても、それらの意味が明らかになるでしょうか？その場合、「点」や「直線」という用語が同時に提示されていなければなりません。そうすれば困難は解消されるでしょう？

この例では、説明は説明すべき用語よりも遥かに難しくなっています。これは常に最も基本的な概念をさらに単純化しようと試みる愚かな行為の結果です。

私たちの数法と数を書き表す方法についての知識は非常に早期に教えられるため、その手法で始まったこと自体がほとんど思い出されていない。したがって、算術の非常に初歩的な部分や計算の簡潔さと洗練さについて考えることは少ない。私たちは十進法による数法を手にしているが、それがどのようにして得られたのかほとんど知らないし、その自然で明白な印象から、この方法は言語と共に生まれたものであり、私たちの言語の一端であると考える。そして、そのためには教育的な指導に多く依存していないと結論付ける。

よく言われているように、もし地球上すべての人々が同じ言語を話していたら、何かの名前はただその対象を表すために選ばれた記号であるという考えではなく、その音が何かと実質的に関連していると考えるだろう。そして、どの他の音でも同様に良いと主張する人に対して笑い、あるいは迫害することも考えるだろう。

このような誤りには陥らないのは、私たちは「馬」を「horse」と呼ぶ音と、「equus」という音で呼ばれるようにローマ人が区別したことを知っているからである。しかし、私たちの数法について同じような誤りをしていると言える。現在では、すべての文明国家がこの方法を受け入れているため、かつてほとんど全人類が、そして今でも原始人たちはその他の方法を使用していたという事実を考慮に入れることなく、私たちは考えている。

以下に述べる考察は、このような状況を正しく捉え、算術に関する私たちの概念においてまだ神秘主義的な要素が残っていることを示すかもしれません。数に対する神秘主義は古今を通じて広く蔓延していたからです。

九つの符号で前九個数字を表し、零の符号も使用します。十を表すために新しい符号を使わず、他の符号を組み合わせます。なぜ十を選んだのか——九や八、あるいは十一を選ばなかったのか？私たちは手数法を考えるとき、この理由が明らかになります。数えるときに、我々は一回に一つの指を立てて記憶します。十まで数えた後、再び最初から始めます。これにより、我々は十個ごとにグループ化し、それを「一個」の単位と見なします。同様に、二十もまた「一個」の単位と考えます。このようにして、我々は「百」「千」といった大きな数を形成します。

次の問題は、これらの数を短く便利な方法でどのように表すかです。
\index{記号!一般原則|EtSeq}\%
\index{数の表現|EtSeq}\%

これは、「選択の問題」であることが明らかです。仮に十の位を一つのアクセント、百の位を二つのアクセント、千の位を三つのアクセントで区別するとします。この数を以下のように表すことができます：---
$76' 5'' 4'''$, $6' 75'' 4'''$, $6' 4''' 5'' 7$, $4''' 5'' 6' 7$、そのような方法は全部で24通りあります。しかし、これ以上は必要ありません；数を表す確実な方法は一つ十分です。最も自然な方法は、大小順に並べることで、最大の集合から始めるか最小の集合から始めるかです。したがって $4''' 5'' 6' 7$ や $76' 5'' 4'''$ のいずれかですが、これらの中から最初のものを選びます。

このように数を記述すると、アクセントは不要であることがすぐに明らかになります。単に右から二番目の位置

\chapter{Chapter III.: Elementary Rules of Arithmetic.}

\#\#\# 算術！基本的な則

皆さん、こんにちは。今日は算術の基本的な概念について学び直しましょう。中学校や高校で習った知識を思い出すことで、新たな視点から数学の魅力を感じていただければ嬉しいです。

まず、代数記号 $+$, $-$, $×$, $=$ そして除法を表す記号について考えてみましょう。これらの記号は非常に基本的で、理解するのは意外と簡単です。例えば、$5 + 3$ という式は「5を3に加える」という意味と同じです。しかし、学生が初めて代数を学ぶときには、文字や数値表現、既知量と未知量の区別、そして文字の指数としての数字の使用など、複数の新しい概念を一度に習得しようとすると混乱することがあります。

これらの記号は一度に一つずつ理解していくことが大切です。例えば、加法 $2 + 7 = 9$ や減法 $6 - 4 = 2$ のような式を頻繁に書きましょう。さらに複雑な式も練習してみましょう：
\begin{gather*}
1 + 8 + 4 - 6 = 4 + 2 + 1, \\
2 × 2 + 12 × 12 = 14 × 10 + 2 × 2 × 2.
\end{gather*}

これらの式は、文字を書く練習を通じて記号の意味に慣れるのに役立ちます。計算を行う習慣が形成されるタイミングは、これらの操作記号を使用するのに最も適した時期です。各操作とその記号を常に結びつけることを忘れないようにしましょう。

例えば、$4 × 7$ は $7 × 4$ とは異なる操作を表します。両者の結果はともに28ですが、前者は4を7回加算し、後者は7を4回加算するという意味が異なります。そして $4 × 7 = 7 × 4$ は証明すべき命題であり、仮定されるものではありません。

また、分数の扱いも重要です。例えば $\frac{1}{7} × 4$ と $\frac{4}{7}$ は異なる操作を表す記号ですが、その結果は同じであることを、分数の扱いを通じて示します。

初学者が練習に選ぶ例題は、単純で非常に大きな数を含まないべきです。心の力は一度に二つ以上のことに向かうことができません：もし使用される数値の複雑さが学生の全注意力を必要とするなら、彼はその規則の原理を見逃す可能性があります。

初学者が学業を開始する段階では、教官によって説明された原理はすぐに学生によって理解されることではない。彼はそれが一般化されるわけではなく、教えられるまでにはならない；彼はそれを学ぶために教えられるべきである；そして特定の事実がすべての数に対して成り立つことを学ぶためには、それがいくつかの数に対して成り立つことが示され、それらのいくつかの数に他の数を代入し、同じ証明を使用することが必要です。彼自身でこれをできるまでには理解しておらず、これは規則が説明され、小さな数を使って自分で試すことでしか達成できない。

例えば、最初は $2 + 3 = 5$ や $4 - 1 = 3$ のような簡単な式から始めましょう。これらの基本的な操作を理解し、自分自身で確認してみることで、規則がどのように機能するかを体験できます。

数学の真理は、一度に一つずつ学び、理解していくことで得られます。初めは少し難しいかもしれませんが、少しずつ練習を重ねていくと自然と理解できるようになります。

次回は、算術における加減法や乗除法について詳しく見てみましょう。これらの基本的な演算がどのようにして他の計算に結びついているか、そしてその原理がどのように機能するかを深く掘り下げてみましょう。皆さん、一緒に学んでいきましょう！

\chapter{Chapter IV.: Arithmetical Fractions.}

\#\#\# Chapter IV.: Arithmetical Fractions

\#\#\#\# 分数の基本：適切な理解から始める

学生が四則演算と最大公約数を求めることのルールを完全に習得した後、彼はすぐに分数の部分に進むべきです。この算術の部分は通常、非常に難しいものとして扱われますが、事実は、分数の中で原則や実践において特に難しいものは、最大公約数を見つけ出すルールだけであるということです。

まず、分数とは何かを理解しましょう。例えば、$56$ が $8$ で割られる場合、その過程は次のように書きます：$\frac{56}{8}$。これは $56$ を $8$ 個の等しい部分に分割し、その一つを商と呼びます。この場合、商は $7$ です。

しかし、$57$ も同様に $8$ 個の等しい部分で分割できることが可能です；例えば、$57$ フィートを $8$ 個の等しい部分に分割できますが、その八分の一は正確なフィート数ではなくなります。なぜなら $57$ には正確な $8$ の個数が含まれていないからです。商 $\frac{57}{8}$ には一部のフィートが含まれることになり、この商は分数や破り数と呼ばれます。

もし $57$ を $56$ と $1$ に分割し、それぞれの八分の一を取ると、全体の八分の一を求めることができます。$56$ フィートの八分の一は $7$ フィートです；一方、$1$ フィートの八分の一は分数であり、これを $\frac{1}{8}$ と書き表します。したがって、$\frac{57}{8}$ は $7 + \frac{1}{8}$ となり、通常は $7\frac{1}{8}$ と書きます。

これらの両方の量 $\frac{57}{8}$ と $7\frac{1}{8}$ は分数と呼ばれますが、唯一の違いは、後者の商の中で整数部分が任意の整数未満の部分から分離されていることです。つまり、$7 + \frac{1}{8}$ のように、整数部分と分数部分が明確に区別されています。

分数は二つの方法で考えることができます。例えば、$\frac{5}{8}$ を取り扱います。これは $5$ を $8$ 個の部分に分割し、その中の $1$ つを表します。同じ長さを得るために、$1$ を $8$ 個の部分に分割して $5$ 個を取り、または $\frac{1}{8} × 5$ を計算することもできます。

これを証明するために、下記のように線を引きます。各線は $1$ インチの $\frac{1}{8}$ を表します；$1$ インチを $8$ 個に繰り返し分割して $5$ 個を取ります。

\[
\ArrayCompress[0.8]
\begin{array}{c*{4}{>{\qquad}c}}
\BAR & \BAR & \BAR & \BAR & \BAR \\
\BAR & \BAR & \BAR & \BAR & \BAR \\
\BAR & \BAR & \BAR & \BAR & \BAR \\
\BAR & \BAR & \BAR & \BAR & \BAR \\
\BAR & \BAR & \BAR & \BAR & \BAR \\
\BAR & \BAR & \BAR & \BAR & \BAR \\
\BAR & \BAR & \BAR & \BAR & \BAR \\
\BAR & \BAR & \BAR & \BAR & \BAR
\end{array}
\]

各列には $1$ インチの $\frac{1}{8}$ が $8$ 回繰り返されています；つまり、$1$ インチです。したがって、全部で $5$ インチがあります。なぜなら、$5$ 列があるからです。しかし、$8$ 行があり、各行は $5$ インチの $\frac{1}{8}$ であり、つまり $\frac{5}{8}$ ですが、各行は $\frac{1}{8}$ の $5$ 回繰り返しであるため、$\frac{1}{8} × 5$ です。したがって、$\frac{5}{8} = \frac{1}{8} × 5$；つまり、$5$ インチを $8$ 個の部分に分割して $1$ 個を取り、または $1$ インチを $8$ 個の部分に分割して $5$ 個を取ります。符号 $\frac{5}{8}$ は両方の操作を表すために使われ、結果が同じであることを示します。

分数の最も重要な性質は、分子と分母に同じ数を掛けることで、その分数の値が変化しないことです。すなわち、$\frac{3}{5}$ は $\frac{12}{20}$ と同じであり、それぞれの部分は $12$インチを $20$部に分けるときと、$3$インチを $5$部に分けるときに同じです。また、$1$インチを $20$部に分割し、その中の $12$部を取り、これは $1$インチを $5$部に分割し、その中的の $3$部を取り出すことと等価です。これにはほとんど証明の必要はありません。$20$から $12$を取ることは、$5$から $3$を取ることと同じであり、$12$が含むたびに $20$の中に $5$が含まれるからです。したがって、どの分数もその形は無数に変更でき、値には影響を与えません。例えば、$\frac{1}{2} = \frac{2}{4} = \frac{3}{6} = \frac{4}{8} = \frac{5}{10}$ など；$\frac{2}{7} = \frac{4}{14} = \frac{6}{21} = \frac{8}{28}$ など。

この原理に基づいて、分数の項は任意の数で割ることでもその値に変化はありません。これにより、共通分母にする、分数を最も簡単な形にする、加算や減算を行う際に困難はありません。これらの操作に関するルールはすべて、算術の本には記載されています。

\#\#\#\# 分数の進化と普遍性

現在、私たちは一つの規則について話しています：分数の値が変化しないこと。しかし、この規則はさらに広範な概念を含んでいます。最初に、私たちは特定の数字（例えば1, 2, 3, 20など）から始めました。その後、これらの数字を部分に分割し、その一部を取り扱いました。

ここで重要なのは、最初の規則がより広範な概念の一部であるという事実です。つまり、私たちは最初は特定の数字（1, 2, 3, 20など）について考えましたが、後にこれらの数字を部分に分割し、その一部を取り扱うことで、更なる普遍性と一般的性を得ました。

例えば、最初は $56$ を $8$ 個の等しい部分に分割して $7$ の商を求めた。しかし、これは更なる一般化へと繋がります：任意の数を任意の他の数で割ることでも、その値は変化しません。つまり、$\frac{56}{8} = 7$ と同じように、$\frac{120}{20} = 6$ や $\frac{35}{7} = 5$ のような操作も成り立ちます。

この原理に基づいて、分数の項は任意の数で割ることでもその値に変化はありません。これにより、共通分母にする、分数を最も簡単な形にする、加算や減算を行う際に困難はありません。これらの操作に関するルールはすべて、算術の本には記載されています。

\#\#\#\# 分数の普遍性と名前の拡張

最後に、分数の普遍性について考えてみましょう。最初の規則が特定の数字（1, 2, 3, 20など）について適用されたとき、その名前は狭隘で片面的でした。しかし、更なる研究と反芻を通じて、私たちは同じ操作を異なる数値にも適用することが可能だと理解しました。

例えば、最初に $\frac{5}{8}$ を定義したとき、その名前は特定の数字（5と8）についてのみ適用されました。しかし、更なる研究により、$\frac{10}{16} = \frac{5}{8}$ など、異なる数値でも同じ操作が成り立つことが分かりました。

このように、最初の規則が更なる普遍性をもたらし、名前の拡張を必要としました。分数の定義は、より広範な概念へと進化し、その適用範囲が広がりました。

分数は、私たちの数学的な理解を深め、より広範で美しい世界を展開する鍵となります。最初から最後まで、分数は私たちの思考を導き、新たな視点を提供します。

\chapter{Chapter V.: Decimal Fractions.}

\#\#\# Chapter V.: 小数分数

\#\#\#\# 前言

皆さん、こんにちは！今日は小数分数について学び直しましょう。小数分数は、私たちの日常生活や科学技術で頻繁に使われる重要な概念です。しかし、学生時代に数学で挫折してしまった方々も安心してください。ここでは、中学校・高校レベルの知識を基に、丁寧に解説します。

\#\#\#\# 小数分数とは？

小数分数は、通常の分数と似ていますが、分母が10, 100, 1000など10の累乗である分数です。例えば、$\frac{75}{100}$や$\frac{333}{1000}$などが小数分数にあたります。

\#\#\#\#\# 分数と小数分数の違い

初心者の多くは、通常の分数（ vulgar fraction）と小数分数が全く異なる分野であるように感じることがあります。しかし、実際には、これらの分数を表す方法が異なりるだけです。例えば、$\frac{75}{100}$は小数として$0.75$と表されます。

\#\#\#\#\# 分母のゼロの数

分母が10, 100, 1000などである分数は、その分母に含まれるゼロの数を示すために特別な記号を使います。例えば、$\frac{17334}{1000}$は$17.334$と表されます。この場合、分子から与えられたゼロの数だけ桁を点線で区切ります。

\#\#\#\#\# 分母が不足する場合

しかし、分子に十分な桁がない場合、例えば$\frac{88}{10000}$はどのように表すべきでしょうか？ここでは、分母のゼロの数を示すために分子の最初にゼロを追加します。つまり、$\frac{88}{10000}$は$0.0088$と表されます。

\#\#\#\# 小数分数の表現方法

\#\#\#\#\# 例題：小数分数の表現

例えば、分数$217.3426$を考えてみましょう。この分数は$200 + 10 + 7 + \frac{3}{10} + \frac{4}{100} + \frac{2}{1000} + \frac{6}{10000}$と分解できます。

\#\#\#\#\# 小数点の位置

小数点左側の桁は通常の数として計算され、小数点右側の数字は右に進むにつれて減少し、各桁は小数点一つ近づくと十分の一になる値となります。例えば、$0.68342$は$0.6 + 0.08 + 0.003 + 0.0004 + 0.00002$と分解できます。

\#\#\#\#\# 演習問題

学習者は他の分数について同様の調査を行い、以下の例題のような演習を分数の原理を使って証明し、この新しい分数記法に完全に慣れさせるべきです：

\[
\begin{aligned}
.68342 &= .6 + .08 + .003 + .0004 + .00002, \\
\text{または}\quad
\frac{68342}{100000}
  &= \frac{6}{10} + \frac{8}{100} + \frac{3}{1000}
  + \frac{4}{10000} + \frac{2}{100000}, \\
.00012 &= .0001 + .00002 = \frac{1}{10000} + \frac{2}{100000}.
\end{aligned}
\]

\#\#\#\# 小数分数の四則演算

\#\#\#\#\# 加法と減法

小数分数の加法と減法は、小数点を同じ列に保つことで実現します。これは、通常の分数を共通分母に変換することと同じです。

例えば、$1.5$と$2.125$を考えます。最初の分数を第二の分数の分母に変換すると、$\frac{1500}{1000}$と$\frac{2125}{1000}$となります。分子を足し合わせると、分数の和$\frac{3625}{1000}$（または$3.625$）が得られます。

\#\#\#\#\# 乗法

二つの小数を掛け算すると、その積には両方の小数に含まれるゼロの数だけの桁が含まれます。例えば、$100 × 1000 = 100000$, $10 × 100 = 1000$などです。

\#\#\#\# 結論

今日は小数分数について学び直しました。小数分数は日常生活や科学技術で頻繁に使われる重要な概念です。これからも、数学の奥深さを楽しみながら学んでいきましょう。

皆さん、頑張ってください！

\chapter{Chapter VI.: Algebraical Notation and Principles.}

\#\#\# 代数学の表記法

\#\#\#\# 短縮符号と表記法

私たちが日々使う言葉や記号は、その言葉や記号が持つ意味を簡潔に表現するためのものであり、特に数学においては重要な役割を果たします。例えば、数値を表すために「8」という言葉を使う代わりに、「八」という言葉を使わず、「8」で表すように短縮しています。これは最初のステップです。

さらに、これらの記号がどのように組み合わさるかによって新たな意味を持つようになっています。例えば、「187」は「100を80に加え7に加えた」という意味になります。「5 + 8」は「5に8を加える」という意味を持ちます。

これらの記号を使用することで、論理的な推論を行い、全ての数に対しても通用する真実が発見されました。例えば、「二つの数の合計の半分に差の半分を加えると、大きい方の数になります」。これは具体的な例として、16と10の場合を考えます。それらの合計の半分は13、差の半分は3；13と3を足すと16となり、その二つの数のうち大きい方です。

他の任意の数（27と8、15と19など）でも同じ命題が成り立つことを確認できます。これらの結果を記号を使用して表すと、以下の真実が得られます：

\[
\frac{16 + 10}{2} + \frac{16 - 10}{2} = 16, \\
\frac{27 + 8}{2} + \frac{27 - 8}{2} = 27, \\
\frac{15 + 9}{2} + \frac{15 - 9}{2} = 15.
\]

任意の二つの数を選び、最初の数と第二の数と呼び、最初の数が大きい場合に最初の数と呼びます。このとき以下の式が成り立ちます：

\[
\frac{\text{最初の数} + \text{第二の数}}{2}
+ \frac{\text{最初の数} - \text{第二の数}}{2}
= \text{最初の数}.
\]

このような方法で、全ての数に対しても通用する何かを表現するために、「First No.」や「Second No.」などと異なる数を表す記号を使用し、その後「First No.」や「Second No.」などを任意の数に置き換えることができます。このようにして以下の主張が常に成り立つことを確認できます：

\[
(\text{最初の数} + \text{第二の数}) ×
(\text{最初の数} - \text{第二の数})
= \text{最初の数} × \text{最初の数} - \text{第二の数} × \text{第二の数}.
\]

\#\#\#\# 方程式

任意の文が二つの数や数の集合が互いに等しいことを述べた場合、それを「方程式」と呼びます。例えば、「7 + 5 = 12」は方程式であり、先に述べた文も方程式です。

次に、頻繁に記述する必要がある「First No.」「Second No.」などを避けることができないでしょうか？これはアルファベットの文字を使用してこれらの数を表すことで実現されます。例えば、「First No.」を「x」とし、「Second No.」を「y」とすると、既に述べた二つの主張は以下のようになります：

\[
\frac{x + y}{2} + \frac{x - y}{2} = x, \\
(x + y) × (x - y) = x × x - y × y.
\]

文字を使用することで、すべての数について何かを述べる文を非常に短時間と少量の労力で記述できるようになりました。現在、使用される様々な符号を列挙します。

1. **アルファベットの文字は数を表すために使用され**ます。文字が用いられた場合、その文字に代わる任意の数値を使用するか、またはその文字が表す数値が未知であり、文字がその位置を占めるまで、その文字が理由づけの中で場所を占めることを意味します。

2. **足し算には「+」の符号が使用され**ます。例えば、「x + z」は「x」と「z」を表す数の和です。以下の方程式は十分明らかです：

\[
x + y + z = x + z + y = y + z + x. \\
a = b \text{ のとき } a + c = b + c, \quad
a + c + d = b + c + d, \quad \etc.
\]

3. **減算には「-」の符号が使用され**ます。以下の方程式はその使用を示しています：

\[
x + a - b - c + e
  = x + a + e - b - c \\
  = a - c + e - b + x.
\]

\[
a = b \text{ のとき } a - c = b - c, \quad
a - c + d = b - c + d, \quad \etc.
\]

4. **乗法の記号「×」は算術で用いるが、文字で表された2つの数を乗じた場合**、「a × b」を「ab」と書くことで表現できるので無駄である。また、「a × b × c」は「abc」、そして「a × a × a」は一時的に「aaa」と表すことにする。2つの数を乗じたときには、乗法の記号「×」 を省略すると誤解されるため必要である。つまり、7 × 5 または 35 は 75 と誤認されてしまうからである。しかし、2つの数を乗じる際に点を挟むのが一般的であり、例えば 7 × 5 × 3 は 7·5·3 と書く。しかしながら、この点が小数点と混同されることがあるため、常に小数点を数字の頭に書き込むことでこれを避けることができる。例えば、23461/100 を 234\Dpt 61 と表す。

5. **除法は算術と同じように記述する**：つまり、「\(\dfrac{a}{b}\)」は数「a」が数「b」で割られる事を意味する。

6. **全ての数の集合を式と呼ぶ**：例えば、\(a + b\)、\(a + b - c\)、\(aa + bb - d\) は代数学的式である。

7. **2つの式を乗じる場合、それぞれを括弧で囲み並べることにより**、その積が明確になります。例えば、「(x + y)(z + w)」は「x + y」と「z + w」の積を表します。

これらの記号と方程式の使用方法を理解することで、数学の世界にさらに深く入り込んでいくことができます。ぜひ、これらの知識を使って新たな発見を楽しんでください！

\chapter{Chapter VII.: Elementary Rules of Algebra.}

\#\#\# Chapter VII.: Elementary Rules of Algebra

\#\#\#\# 加法と減法のアルゴリズムについて

学生は、加法や減法のアルゴリズムを学ぶ前に、これまで定めた原理や表記法について非常に詳しく理解しているべきです。それから、それぞれの単純な例を取り上げて、次のように論理的に合計と差を求めることを進めるべきです。

例えば、$c - d$ を $a - b$ に加える必要がある場合を考えましょう。この指示は通常の方法で書くことができます：

\[
\begin{array}{rc}
    & a - b \\
    + & c - d \\
\cline{2-2}
\Strut\text{Add }
\end{array}
\]

または、より適切には $(a - b) + (c - d)$ を求めることです。

もし $c$ を $a$ に加えると、つまり $a + c$ を求めると、多すぎるかもしれません。最初は、増加すべきは $a - b$ であり、$b$ が減らされるべきであるからです；したがって、$a + c - b$ になります。しかし、それでもまだ多すぎます。なぜなら、加えるべきは $c - d$ であり、$d$ が減らされるべきであるからです；したがって、$a + c - b - d$ または $a - b + c - d$ になります。

このような数回の論理的な推論からルールを導き出すことができます。そしてその時までに、学生が一瞬でも証明を見失うような複雑な例を選ばないでください。

\#\#\#\# 乗法について

減法についてはすでに実行していますので、いくつかのこの方法の例からルールを導き出すことができます。

まず、$a$ を $c - d$ 倍する方法は次の通りです：
\[
a
\]
を最初に $c$ 回取るか、または $ac$ を求めます。これがあまりにも大きいため、$a$ が $d$ よりも多く取られています。したがって、$ac$ から $d$ 倍の $a$，すなわち $ad$ を引く必要があります。結果は次のようになります：
\[
a(c - d) = ac - ad.
\]
これは算術から検証できます。同じ過程が正しいと示されています；そして、この方法は整数や分数であっても有効です。例えば、$6(14 - 9)$ または $6 \times 5$ は $6 \times 14 - 6 \times 9$，すなわち $84 - 54$ と等しいことがわかります。また、$\frac{2}{3}(\frac{1}{7} - \frac{2}{15})$ または $\frac{2}{3} \times \frac{1}{105}$ は $\frac{2}{3} \times \frac{1}{7} - \frac{2}{3} \times \frac{2}{15}$，すなわち $\frac{2}{21} - \frac{4}{45}$ と等しいことがわかります。同様の論理に基づいて、以下の式も証明できます：
\begin{gather*}
a(b + c - d) = ab + ac - ad, \\
(p + pq - ar)xz = pxz + pqxz - arxz, \\
(a^{2} + 2b^{2})b^{2} \text{または} (aa + 2bb)bb = aabb + 2bbbb = a^{2} b^{2} + 2b^{4}.
\end{gather*}
また、乗算が行われた場合、その過程を逆に進めることで因子を求めることができます。例えば、式
\[
ab - ac + a^{2}, \quad \text{または} \quad ab - ac + aa,
\]
の形成を見ると、各項はすべて $a$ で乗算されています；つまり、これは $(b - c + a)$ を $a$ で乗算したものです。これにより
\[
ac + ad + bc + bd = a(c + d) + b(c + d) = (a + b)(c + d),
\]
および
\[
a^{2} - ab^{2} + 2abc - dc + 3a = a(a - b^{2} + 3) + c(2ba - d).
\]
を理解することが容易になります。算術では、整数や分数であっても順序を変えても積は同じであることが証明されています。すなわち $6 \times 4 \times 3 = 3 \times 6 \times 4 = 4 \times 6 \times 3$ 等が成り立ちますし、また一部の乗算を先に進め、その部分積で元の因子を置き換えることもできます。例えば $3 \times 4 \times 5 \times 6$ は $12 \times 5 \times 6$，または $15 \times 4 \times 6$，または $90 \times 4$ と表せます。これらのルールから、複雑な単項式を乗算し、可能な最も簡単な形でその積を表現することができます。例えば
\[
6a^{3}b^{4}c, \text{これは } 6aaa bbbb c,
\]
を $12a^{2}b^{3}c^{3}d$，すなわち $12aa bbb ccc d$ で乗算する場合、積は
\[
6aaa bbbb c \times 12aa bbb ccc d
\]
と書き、この乗算は次の順序で行われます：
\[
6 \times 12 aaaaa bbbbbbb cccc d,
\]
これは $72a^{5}b^{7}c^{4}d$ と表されます。このような種類の例をいくつか見ることで、代数の教科書で通常与えられるような量の乗算のルールが確立します。

\#\#\#\# 代数的な公式

すべての数について成り立つ代数的な公式において、その中の一つの文字に対して任意の代数的表現を代入することが可能であり、その文字が公式中にどの場所でも現れる限り注意深く代入する必要があります。したがって、次の公式：
\[
a^{2} - b^{2} = (a + b)(a - b)
\]
から、$a$ に対して代入することにより以下の式を導けます。この公式が常に成り立つ場合、$a$ が $p + q$ に等しいときも成り立ち、つまり公式中の $a$ のどの場所でも $p + q$ を代入すればよいです。したがって，
\[
(p + q)^{2} - b^{2} = (p + q + b)(p + q - b).
\]

同様に，
\begin{gather*}
(b + m)^{2} - b^{2} = (2b + m)m, \\
(x + y)^{2} - (x - y)^{2}
  = (x + y + x - y)(x + y - \overline{x - y})
  = 4xy,
\end{gather*}
と続きます。

既に公式
\[
(p - q)a = ap - aq.
\]
を証明しました。

この式の $a$ を $r - s$ に置き換えると、次のようになります：
\[
(p - q)(r - s) = (r - s)p - (r -
\]

以上で、加法、減法、乗法、そして代数的な公式について詳しく説明しました。これらのルールは、中学校や高校の数学の知識をベースにしているので、大人の読者でもスッキリと理解できるはずです。ぜひ一度試してみてください！

\chapter{Chapter VIII.: Equations of the First Degree.}

\#\#\# Chapter VIII.: 方程式の第一度

我々は既に方程式の定義を述べており、様々な種類の多くの方程式に出会いましたが、これら全てが特定の数によって表される文字に依存しない性質を持っていました。つまり、どの文字に対しても同じ数値を代入しても等式は常に成り立つものでした。例えば、

\[
\frac{a^{2} - 1}{a + 1} = a - 1
\]

この代数的な命題の真偽は、$a$が$1$, $2$, $2\frac{1}{2}$などどのような数値を表しても同じです。実際、方程式の右辺は左辺で示された操作の結果を表しています。左辺では$a^{2} - 1$を$a + 1$で割ることを求められています。右辺はその割り算の結果を示しています。このような性質を持つ方程式は、実際には両辺が同じ数値を異なる表現方法で記述しているため、「恒等式」と呼ばれます。

以下にいくつかの恒等式を示します：

\[
a + a = 2a, \quad
7a - 3a + b = 4a - 3b + 4b, \quad \text{および} \quad
\frac{a}{b} × b = a.
\]

最後の項の末尾にある全ての公式は恒等式の例です。これらの一つもまた、その中の文字に与えられるすべての値に対して常に真であり、ただし一方の部分で文字が表すものが他方の部分でも同じである限りです。

もし方程式 $a + 1 = 8$ を考えるなら、これはその代数的な量に対してどのような値でも真となるわけではありません。明らかにこの方程式を満たす唯一の数は $a = 7$ であり、他の仮定は矛盾を生じます。これは新しい種類の方程式で、特定の状況下でしか存在しないものです。その特定の状況は方程式自体から求めることができます。問題の解法が導くような方程式は、後ほど示すように常にこのようなものであり、代数学の基本ではこの種類の方程式が最も重要です。これらの方程式を恒等式と区別するために、「条件方程式」と呼ばれます。これはその中に含まれる文字がどのような数でも表すことができない限り真になりませんし、その存在自体は文字が満たすべき条件を作ります。条件方程式の解法とは、方程式が真となるために文字が何であるべきかを見つける過程です。このようなすべての解法は論理的なプロセスであり、方程式の真偽を仮定し、自明な手順で進行し、算術や代数学の一般的なルールを使用します。

問題の解法において条件方程式に至るまでの例を見た後、その解法について再び扱います。条件方程式では、その方程式によって決定される量は通常最後のアルファベットの一つの文字で表されますし、他のすべては最初のアルファベットのいくつかで表されます。この区別は初心者にとって必要ですが、時間とともに彼はそれを捨てて、どの文字が問題の条件によって既知または未知の数を表すかを考えるようになるべきです。

以下の問題を代数方程式に簡約する一般的なルールは与えられません。問題は、ある数の性質が言葉で表現され、その数を見つけ出すためにその性質が必要です。この性質は最も便利な方法で方程式として書き下されるべきです。以下に、この例を示す問題の簡約化された方程式を挙げます：

\#\#\#\# 1. 56を加えたときに200より2倍少ない数が得られる数は何ですか？

求めたい数を$x$とします。

すると、$x + 56 = 200 - 2x$となります。

56, 200, 2の代わりに任意の他の与えられた数$a$, $b$, $c$を使用しても同じ方法で、$x$を求め出す方程式は
\[
x + a = b - cx.
\]
となります。

\#\#\#\# 2. 同じ場所から出発した2人の信使がいます。後ろの信使者は1時間に3マイルを、前の信使者は2マイルを進みます。前の方の信使者が出発して4時間後に後ろの信使者が追いつくまでに何時間かかりますか？

追いつくまでの時間を$x$とします。この出来事が起こるとき、前の方の信使者は$x + 4$時間走り、$(x + 4)2$マイル、すなわち$2x + 8$マイルを進みます。後ろの信使者は$x$時間走り、$3x$マイルを進みます。そして、追いつくときには両者とも同じ距離を進んでいるため，
\[
3x = 2x + 8.
\]
となります。

前の方の信使者が1時間に$a$マイルを、後ろの信使者が$b$マイルを、さらに$c$時間を経過した後に後ろの信使者が追いつくとき、方程式は
\[
bx = ax + ac.
\]

\#\#\#\# 3. 四人、A、B、C、Dが船を造り、その費用は£2607であった。BはAの二倍の金額を支払い、CはAとBの合計と同じ金額を支払い、DはCとBの合計と同じ金額を支払った。それぞれがどのくらい支払ったかを求めよ。

\begin{align*}
\text{もし } \&A \text{ が } x \text{ ポンドを支払ったとするなら,} \\
\text{すると } \&B \text{ は } 2x \text{ ポンドを支払った}\dots\Add{,} \\
              \&C \text{ は } x + 2x \text{ もしくは } 3x \text{ ポンドを支払った}\dots\Add{,} \\
              \&D \text{ は } 2x + 3x \text{ もしくは } 5x \text{ ポンドを支払った}\dots\Add{.}
\end{align*}

\#\#\#\# 4. 二つの栓があり、一つ目は12時間で水槽を満たし、もう一つは15時間で満たす。両方を開けて同時に流した場合、どのくらい時間がかかるか。

$x$時間を経過するまでに水槽が満たされるとする。すると、一つ目の栓は12時間で水槽を満たすので、1時間では$\frac{1}{12}$を満たし、2時間では$\frac{2}{12}$を満たし、$x$時間では$\frac{x}{12}$を満たす。同様に、二つ目の栓も$x$時間で$\frac{x}{15}$を満たす。両方の栓が水槽を完全に満たしたとき、これらの分数の合計は1 entiretyを表すので,
\[
\frac{x}{12} + \frac{x}{15} = 1.
\]
二つの栓で水槽を満たす時間はそれぞれ$a$と$b$とするとき、方程式は
\[
\frac{x}{a} + \frac{y}{b} = 1.
\]

\#\#\#\# 5. ある人が8ヤードの布を£32先令で購入し、一部は9先令1ヤード、残りは7先令1ヤードで買った。それぞれどのくらいの量を買ったか。

$x$を7先令1ヤードで買ったヤード数とする。すると$7x$はその価格である。また、8−xは9先令1ヤードで買ったヤード数であり、$(8 - x)9$もしくは$72 - 9x$がその価格である。これら二つの合計、すなわち$7x + 72 - 9x$が全体の価格£32先令（62先令）であることから,
\[
7x + 72 - 9x = 62.
\]

これらの例は、問題を方程式に簡約する方法が示されるのに十分である。未知数を表す文字を仮定し、

\#\#\# 解説講義

この章では、方程式の第一度について深く学び直します。まず、恒等式と条件方程式の違いについて理解しましょう。

\#\#\#\# 恒等式
恒等式とは、特定の数値に依存しない方程式のことです。例えば、

\[
a + a = 2a, \quad 7a - 3a + b = 4a - 3b + 4b, \quad \text{および} \quad \frac{a}{b} × b = a
\]

これらの式は、どの数値を代入しても常に成り立ちます。例えば、

\[
\frac{2^2 - 1}{2 + 1} = 2 - 1
\]

この等式は、$2$以外の任意の数値を代入しても成立します。

\#\#\#\# 条件方程式
一方、条件方程式は特定の状況下でしか成り立たない方程式です。例えば、

\[
a + 1 = 8
\]

この方程式は、$a = 7$ のときにのみ成り立ちます。他の数値を代入すると矛盾が生じます。

\#\#\#\# 方程式の解法
条件方程式を解くためには、未知数を特定の文字で表し、その文字がどのような数値になるべきかを見つける必要があります。例えば、

1. 56を加えたときに200より2倍少ない数が得られる数は何ですか？

この問題は方程式

\[
x + 56 = 200 - 2x
\]

に簡約されます。これを解くと、

\[
3x = 144 \implies x = 48
\]

となります。

2. 同じ場所から出発した2人の信使がいます。後ろの信使者は1時間に3マイルを、前の信使者は2マイルを進みます。前の方の信使者が出発して4時間後に後ろの信使者が追いつくまでに何時間かかりますか？

この問題は方程式

\[
3x = 2(x + 4) \implies x = 8
\]

に簡約されます。

3. 四人、A、B、C、Dが船を造り、その費用は£2607であった。BはAの二倍の金額を支払い、CはAとBの合計と同じ金額を支払い、DはCとBの合計と同じ金額を支払った。それぞれがどのくらい支払ったかを求めよ。

この問題は方程式

\[
x + 2x + (x + 2x) + (x + 2x + x + 2x) = 2607
\]

に簡約されます。これを解くと、

\[
10x = 2607 \implies x = 260.7
\]

となります。

4. 二つの栓があり、一つ目は12時間で水槽を満たし、もう一つは15時間で満たす。両方を開けて同時に流した場合、どのくらい時間がかかるか。

この問題は方程式

\[
\frac{x}{12} + \frac{x}{15} = 1
\]

に簡約されます。これを解くと、

\[
x = 60
\]

となります。

5. ある人が8ヤードの布を£32先令で購入し、一部は9先令1ヤード、残りは7先令1ヤードで買った。それぞれどのくらいの量を買ったか。

この問題は方程式

\[
7x + (8 - x)9 = 62
\]

に簡約されます。これを解くと、

\[
x = 5 \implies 7x = 35, \quad 8 - x = 3
\]

となります。

これらの問題は、方程式の第一度を理解し、解法を学ぶのに十分な例です。未知数を表す文字を仮定し、その文字がどのような数値になるべきかを見つけることで、問題を解決することができます。この過程は論理的なプロセスであり、算術や代数学の基本的なルールを使用します。

\#\#\# まとめ

方程式の第一度について学び直すことで、我々は未知数を見つけ出すための方法を再確認し、問題解決のスキルを高めることができます。恒等式と条件方程式の違いを理解し、具体的な例を通じて解法を練習することで、数学の魅力を深く感じることができるでしょう。

ご自身の経験や知識を活かしながら、これらの概念を学び直すことで、

\chapter{Chapter IX.: On the Negative Sign, etc.}

\#\#\# Chapter IX.: On the Negative Sign, etc.

\#\#\#\# 負の符号について

数学において、負の符号はしばしば混乱を招きます。しかし、この章では、負の符号に関する基本的な概念と操作方法を丁寧に説明します。

まず、8が5より3大きいという意味で述べたい場合、式は$8 - 5 = 3$となります。また、$a$が$b$より$cだけ大きい」という意味で述べたい場合、式は$a - b = c$を使います。値が既知の数を他の同等に知られた数から引く限り、大きい方から小さい方を引こうとする心配はありません。例えば$3 - 8$と書くようなことはありません。

しかし、文字を含む調査を進めているときには、時として注意不足や、どの量が大きいかの知識不足、または問題の一部の条件に対する誤解から、引き算で数を入れ替えることがあります。例えば、二つの量の大きい方$b$が$a$よりも大きい場合に、$a - b$と書く代わりに$b - a$と書いてしまうことがあります。

\#\#\#\# 二つの量の和について

もし二つの量の和について述べたとしても、それは問題ありません。なぜなら$a + b$と$b + a$は同じだからです。しかし、$a - b$と$b - a$は異なる場合があります。例えば、$8 - 3$は理解できます。3から8を引くことは不可能で、3より少ない数を3から引かなければならないということが不合理です。もし問題の答えとして$3 - 8$が得られた場合、それは問題そのものに何らかの不合理性があるか、またはそれを方程式に置く方法に問題があったことを示すでしょう。

\#\#\#\# 負の符号の重要性

しかし、このような答えは現実になることがありますので、学生はどのような誤りがこれにつながるのかを知るべきであり、それによって調査過程にどのように影響するかを理解しなければなりません。我々は初心者に対して、これらの量についての概念を形成する際には経験のみを指南役とし、つまり多くの結果の観察から規則を導き出することを推奨します。理論からは何も学ばないようにしましょう。

負の符号に関する困難性は、彼が混乱することになるような方法で最もよく説明されます：おそらく彼はその困難性自体を見ることはできません；数学の学生にとって常に容易な信頼は欠点であり、彼らを反論から始めさせることこそが大きな勝利です。この誤りが引き起こす効果について、これまでに到達したすべての調査において引数を示すために考察し、経験に基づく規則を導き出すことにします。学生はこれらの規則を記憶しておくべきで、これらは経験から得られたものであり、抽象的な論理からは得られないものです。

\#\#\#\# 代数学の拡張

我々は初心者に対して、これらの量についての概念を形成する際には経験のみを指南役とし、つまり多くの結果の観察から規則を導き出ることを推奨します。理論からは何も学ばないようにしましょう。

負の符号に関する困難性は、彼が混乱することになるような方法で最もよく説明されます：おそらく彼はその困難性自体を見ることはできません；数学の学生にとって常に容易な信頼は欠点であり、彼らを反論から始めさせることこそが大きな勝利です。この誤りが引き起こす効果について、これまでに到達したすべての調査において引数を示すために考察し、経験に基づく規則を導き出すことにします。学生はこれらの規則を記憶しておくべきで、これらは経験から得られたものであり、抽象的な論理からは得られないものです。

\#\#\#\# 代数的形式の変更

表記中の各項の符号が変更されたとき、その式は形式を変更したとされます。例えば $+a - b$ と $+b - a$ は異なる形式であり、もし $a$ が $b$ より大きい場合、前者が正しい形式で後者が誤った形式です。

ある量として $3 - 8$ を別の形で表すルールを拡張します。例えば $+3 - 8$ に他の数を接続すると、$56 + 3 - 8$ とします。これは $56 + 3 - (3 + 5)$ または $56 + 3 - 3 - 5$ または $56 - 5$ と書き換えることができます。したがって、$+3 - 8$ に接続された任意の数は $-5$ に接続された同じ数と同じであり、これから $+3 - 8$ や $3 - 8$ は $-5$ と同じであると述べることができます。つまり $3 - 8 = -5$ です。

\#\#\#\# 方程式の変形

これは方程式 $8 - 3 = 5$ の別の書き方であり、$8$ が $5$ より $3$ 大きいことを等しく示しています。同様に $a - b = -c$ は $b$ が $a$ より $c$ の量で大きいことを示します。

もし $a$ が何もないとすると、この方程式は $-b = -c$ となり、これは $b = c$ を示すため、方程式 $a - b = -c$ を正しい形式に書き換えると $b - a = c$ になり、もし $a = 0$ の場合、$b = c$ となります。これで以下の方程式を理解することができます：
\begin{gather*}
a - b + c - d = -e, \quad \text{または} \quad b + d - a - c = e, \\
2ab - a^{2} - b^{2} = -d - e, \quad \text{または} \quad a^{2} + b^{2} - 2ab = d + e.
\end{gather*}

\#\#\#\# 方程式の操作

我々は、$a - b = -c$ のような方程式に対して何らかの操作を行う前に、その操作方法を確認する必要があります。正しい形の方程式は $b - a = c$ です。これにより、

\[
d + b - a = d + c,
\]

または

\[
d - (a - b) = d + c.
\]

ここで $a - b$ を表す符号 $-c$ に書き換えて、$d - (-c) = d + c$ とします。ここでは $a - b$ に対して操作を行いましたが、これは $a$ が $b$ より小さいため量ではないからです。しかし、現在の目的は、共通のルールをこのような表現に適用することで結果を見極め、実際の形で示すことです。

最初の辺 $d - (-c)$ には意味を付与することができませんが、右辺はその実際の形 $d + c$ を与えます。この表現の意味は、$a - b$ という量であると誤って考えたものの、実際には $a < b$ より量ではないものに対して代数的な減算ルールを適用することで、$b - a$ の実際の量を $d$ に加える結果と同じを得ることです。$d - (-c)$ の形を使用するのは、その実際の形 $d + c$ との関連性を見失うことなく使用できることと、過程の終了時に誤差を正しく修正できるルールを確立できることからです。

\#\#\#\# 結論

我々が採用するルールは、多くの例によって確立されるもので、二つの類似した符号が一緒に現れる場所では、実形式の対応部分には正の符号が、二つの異なる符号が一緒に現れる場所では、実形式には負の符号があるというものです。従って、$d - (-c)$ の実形式は $d + c$ です。

また、方程式 $a - b = -c$ の実形式 $b - a = c$ を取ります。すると $d - (b - a) = d - c$ となり、または $d - b + a = d - c$ となり、または $d + c$ となります。

このように、負の符号に関する理解を深めることで、代数学はより広範に展開し、多くの質問に対する解決策を提供することができます。

\chapter{Chapter X.: Equations of the Second Degree.}

\#\#\# Chapter X.: 方程式の第二度

\#\#\#\# 1. 平方根の概念とその逆操作

数学におけるすべての操作は、その効果が完全に逆となる別の操作と結びついています。例えば、加法と減法、乗法と除法は逆の操作です。つまり、
\[ a + b - b = a \]
であり、
\[ \frac{ab}{b} = a \]

さらに、指数の上げ下げに関連する逆の操作としてルートの抽出があります。根という用語は次のように説明されます：$a^2$ や $a^3$ を $a$ の平方や立方と呼びますから、$a$ は $a^2$ の平方根と呼ばれます。例えば、169が13の平方であるなら、13は169の平方根です。

以下にこの表現を示す表があります：
\[
\begin{align*}
& a は \( a^2 \) の平方根と呼ばれ、記号で \(\sqrt{a^2}\) で表されます, \\
& a は \( a^3 \) の立方根と呼ばれ、記号で \(\sqrt{a^3}\) で表されます, \\
& a は \( a^4 \) の四次根と呼ばれ、記号で \(\sqrt{a^4}\) で表されます, \\
& a は \( a^5 \) の五次根と呼ばれ、記号で \(\sqrt{a^5}\) で表されます, \\
& その他の例も同様です。
\end{align*}
\]

\#\#\#\# 2. 二次方程式の解法

通常の方法は、一次方程式の解法を学んだ後で、任意の代数的量の平方根をすぐに教えることです。しかし、二次方程式の解法を学ぶまでにはこのルールを省略することをお勧めします。

平方根が取り出せる数は比較的少ないことに注意してください。例えば、7は任意の整数または分数自身との乗積で作ることができません。最初の主張は明らかですが、第二の主張は初心者には容易に証明できません。しかし、多くの例を取ることによって、真の分数（分子が分母で測れない分数）は自乗すると整数にならないという事実を確認できます。

例えば、$\frac{4}{2} \times \frac{4}{3}$ または $\frac{16}{9}$ は整数ではありません。したがって、7は整数の平方でも分数の平方でもありませんし、厳密には平方根を持たない数です。しかし、7に非常に近い分数は存在し、その平方根も存在します。この近さは任意の程度まで持続できます。

別の例として、1.4142を自乗すると積は1.99996164となり、これは2から非常に小さな分数0.00003836だけ異なります。したがって、1.4142の平方は2に非常に近い値であり、さらに近い平方を持つ分数も見つけることができます。

次に、次の問題を考えましょう。ある人が2シャリングで布を一定数のヤード分購入し、そのヤード数は彼が1ヤードにつき支払うシャリング数と一致します。彼はどの程度のヤードを購入したでしょうか？この数を $x$ とすると、$\frac{2}{x}$ は1ヤード当たりの価格であり、$x = \frac{2}{x}$ より $x^2 = 2$ となります。これは我々が述べた通り不可能なことですが、つまり、条件を満たす正確なヤード数やその部分は存在しません。しかし、1.4142ヤードはほぼ適切であり、1.4142136ヤードはさらに近い値です。

実際の問題が提案された場合でも、十分に近い解を見つけることは容易です。したがって、平方根を抽出できない解を持つ問題も、平方根への近似によって有用になります。

\#\#\#\# 3. 二次方程式の一般形

方程式の二つ目の次数、一般的に二次方程式と呼ばれる方程式は、未知数の平方が含まれる方程式です。例えば：
\[ x^2 - 3 = 4x^2 - 15 \]
または
\[ x^2 + 3x = 2x^2 - x - 1 \]

これらの方程式を項の移動により常に次の形のいずれかに簡約できます：
\[
\begin{align*}
& ax^2 + b = 0, \\
& ax^2 - b = 0, \\
& ax^2 + bx + c = 0, \\
& ax^2 - bx + c = 0, \\
& ax^2 + bx - c = 0, \\
& ax^2 - bx - c.
\end{align*}
\]

例えば、上記の二つの方程式はそれぞれ \(3x^2 - 12 = 0\)，\(x^2 - 4x - 1 = 0\) の形を取る。これらの方程式に対処するためにはまず、方程式 \(x^2 = a^2\) を取り扱います。

この方程式は \(x^2 - a^2 = 0\)，または \((x + a)(x - a) = 0\) と等しい。二つ以上の量の積がゼロとなるためには、その中の一つがゼロであることが十分であり、したがって以下の方程式から \(x\) の値を導くことができます：
\[
\begin{align*}
x - a &= 0, \\
\text{または}\quad
x + a &= 0.
\end{align*}
\]

これらはそれぞれ \(x = a\)，\(x = -a\) を与えます。

これを明瞭にするために、以下の方程式から \(x\) を求めましょう：
\[
(3x + a)(a^3 + x^3) = (x^2 + ax)(a^2 + ax + 2x^2).
\]

この方程式を展開し、すべての項を一方に移動すると次のようになります：
\[
x^4 - 2a^2 x^2 + a^4 + 2a^3 x - 2ax^3 = 0,
\]
または
\[
(x^2 - a^2)^2 - 2ax(x^2 - a^2) = 0,
\]
または
\[
(x^2 - a^2)(x^2 - 2ax - a^2) = 0.
\]

最後の方程式は \(x^2 - a^2 = 0\)，または \(x^2 = a^2\) のとき成り立つ。これは \(x = +a\)，または \(x = -a\) のときに成り立つ。

元の方程式に \(+a\) を代入すると結果は \(4a \times 2a^3 = 2a^2 \times 4a^2\) となり、\(-a\) を代入すると結果は \(0 = 0\) となる。これは \(+a\) と \(-a\) が両方とも \(x\) の正しい値であることを示す。

以上で二次方程式の基本的な解法を理解しました。これから先、これらの知識を使ってより複雑な問題にも挑戦してみましょう。

\chapter{Chapter XI.: On Roots in General, and Logarithms.}

\#\#\# 1. 指数と根の関係性を再確認する

まず、指数と根についての基本的な定義を思い出しましょう。これまでに学んだのは、ある数 \( a \) の平方、立方、四次冪などはそれぞれ以下のように表されます：

\[
\begin{array}{*{2}{l}>{\quad}lr}
a の平方 & x = aa  & a の平方根 & xx = a \\
a の立方 & x = aaa & a の立方根 & xxx = a \\
a の四次冪 & x = aaaa  & a の四次根 & xxxx = a \\
a の五次冪 & x = aaaaa & a の五次根 & xxxxx = a
\end{array}
\]

これらの定義は、それぞれの数 \( a \) を何回掛ければその値になるかを表しています。例えば、\( 4 \) の平方は \( 16 \)，そして \( 16 \) の平方根は \( 4 \) です。

\#\#\# 2. 冪と根の一般的な記号

これらの定義は、冪（べき）と根（ルート）を表す一般的な記号で書くことができます：

\[
\begin{gather*}
\text{冪}\quad
a^{2}, a^{3}, a^{4}, a^{5} \dots, a^{m}, a^{m+n}, \quad\etc. \\
\text{根}\quad
\sqrt{a}, \sqrt{a}, \sqrt{a}, \sqrt{a}, \sqrt[m]{a}, \sqrt[m+n]{a}, \quad\etc.
\end{gather*}
\]

これらの記号は、それぞれの数 \( a \) の冪や根を表しています。例えば、\( 16 \) の平方根は \( \sqrt{16} = 4 \)，立方根は \( \sqrt{16} \) です。

\#\#\# 3. 冪と根の関連性

これまで冪と根を表す方法には、直接的な関連性がありませんでした。しかし、規則の拡張により、冪と根の関係性が明らかになります。具体的には、次の等式が成り立ちます：

\[
(\sqrt[n]{a})^n = a
\]

例えば、\( \sqrt{16} = 4 \) であれば、\( (4)^2 = 16 \) です。

\#\#\# 4. 冪の法則と根の法則

任意の項を指数 \( n \) で累乗する一般的な法則は、係数を \( n \) 乗し、各文字の指数を \( n \) 倍することです。例えば：

\[
(a^p b^q c^r)^n = a^{np} b^{nq} c^{nr}
\]

これは、任意の単純な項の根を求めることにも適用されます。その対応する累乗がどのように求められるかに従います。法則は、係数の要求される根を引き出し、各文字の指数を根の指数で割ります。

例えば、\( 16a^{12} b^8 c^4 \) の四乗根を求めます：

\[
\sqrt{16a^{12} b^8 c^4}
\]

この表現は以下のようになります：

\[
2a^3 b^2 c × 2a^3 b^2 c × 2a^3 b^2 c × 2a^3 b^2 c
\]

または、

\[
(2a^3 b^2 c)^4 = a^{12} b^8 c^4
\]

したがって、四乗根は \( 2a^3 b^2 c \) です。

\#\#\# 5. 積の根

任意の積，例えば \( AB \) の根を求めるには、その各因子の根を求め、それらを乗じればよい。例えば：

\[
\sqrt{AB} = \sqrt{A} \sqrt{B}
\]

なぜなら、\(\sqrt{A} \sqrt{B}\) を立方すると、

\[
\sqrt{A} \sqrt{B} × \sqrt{A} \sqrt{B} × \sqrt{A} \sqrt{B}
\]

または

\[
\sqrt{A} \sqrt{A} \sqrt{A} × \sqrt{B} \sqrt{B} \sqrt{B}
\]

または

\[
AB
\]

同様にして、一般に、

\[
\sqrt[n]{ABC} = \sqrt[n]{A} \sqrt[n]{B} \sqrt[n]{C}
\]

が証明できます。

\#\#\# 6. 最高次根の分解

任意の式のどの根も最も単純な表現法は、その式を最高次根が取りうる因子とそれ以外の因子に分けることです。例えば、

\[
\sqrt{16a^4 b^7 c}
\]

を求めると、\( 16 \) は \( 8 × 2 \)，\( a^4 \) は \( a^3 × a \)，\( b^7 \) は \( b^6 × b \) と分解できます。したがって、式は \( 8a^3 b^6 × 2abc \) となり、各因子の立方根を求めてから乗じることで、

\[
2ab^2 \sqrt{2abc}
\]

を得ます。

第二の因子には、これまで使用してきた記号で表現できる立方根がなく、しかし \( a \)，\( b \)，および \( c \) が表す数値が分かれば、\(\sqrt{2abc}\)は正確に求められるか、またはそれが不可能な場合は近似によって求めることができます。

\#\#\# 7. 冪の法則と根の法則の拡張

任意の冪の冪についても、一般的な法則があります。つまり、

\[
(a^m)^n = (a^n)^m = a^{mn}
\]

これは、例えば \( (a^2)^4 \) は \( a^{2×4} = a^8 \)，そして \( (a^4)^2 \) も \( a^{4×2} = a^8 \) です。

同様に、

\[
\sqrt{\sqrt{a}} = \sqrt{a}
\]

これは、なぜなら \( a = (\sqrt{a})^6 \)，つまり平方根の立方根は六乗根であるからです。

\#\#\# 感想

これらの法則と定義は、数学の美しい世界を再発見するための鍵となります。冪と根の関係性を深く理解することで、より複雑な問題にも対処できるようになります。ぜひ、これらの法則を心に刻み、日常生活や他の学問分野でも活用してみてください。

\chapter{Chapter XII.: On the Study of Algebra.}

\#\#\# Chapter XII.: On the Study of Algebra

本章では、代数の学習方法についていくつかのアドバイスを提供します。この主題には無限の範囲と応用例があり、どの名前で呼ばれようと、その基礎となる概念が幾何学的であれ如何いかんにせよ、機械的な手続きは代数的であり、既に扱われた主題の前の規則や発展形の継続的な応用を提供します。これは三角法、代数を幾何学に適用する場合、微分 calculus または流動性など、同様の状況が存在します。

\#\#\#\# 1. 理解と確認

任意の代数的著作物を読む際には、まずその中で示される異なる手続きについて完全な理解を得ることが重要です。また、これらの手続きがどのように相互に関連しているかにも注意を払う必要があります。これは単なる読書では達成できません。数学的な作品においては、学生が完全に理解したと言えるまで、心の中で各手続きを埋め尽くす方法で埋めることが不可能です。多くの結果は詳細が省略されており、加算、乗算、平方根の抽出などがあります。これらの結果は学生が信用してはならず、自分の筆を使って計算する必要があります。代数的手続きに取り組んでいる間は手から離さないようにしましょう。

我々が推奨する方法は、各研究全体の象徴的な部分を書き出し、述べられた部分を埋め、異なるステップの関連性を説明するために絶対的に必要な程度の文言的解釈だけを追加することです。これは通常、書く行為が強制する慎重な考慮よりも多くの説明が与えられるとは限りません。このような作業は試みたことがない人には驚くべき労力に思えるかもしれませんが、私たちはこれがあまりにも短い方法であると確信しています。なぜなら、書き込む行為が我々に強制する慎重な考慮は、単なる読書だけで新しい言語で表現された新しい論理を理解しようとすると発生する混乱や困難を避けるのに役立つからです。

この方法で進める中で何か難しさを感じた場合、その部分を完全に書き出してみることをお勧めします。しばしば、混乱の原因となった誤解はその試練を受けないと明らかになりません。それでもまだ不明な点がある場合は、注意深く再考しても解決しない場合、学生はその点についてメモを取り、その後の進行の中で何か類似性のあるものを発見したときにそれを振り返ることで進むべきです。これにより以前不理解だったものが解読できることがよくあります。

純粋な幾何学的な論理では、全体の議論の詳細を書き込む必要は少ないでしょう。なぜなら象徴的な言語が速く理解でき、主題は操作の機械性に大きく依存していないからです。しかし代数的手続きにおいては、そのどの点にもより多くのことがかかっているため、指導者はこの点について特に強調すべきです。

\#\#\#\# 2. 新たな法則や手続き

II. 新たな法則や手続きに到達した際、学生は自分自身がその法則や手続きを理解し、適用できるかを確認するために十分な数の例題を解くべきです。しかし、多くの例題は書籍において形式的に与えられていないため、難しさを感じるかもしれません。それでも、彼は自分で例題を選んで、以前の知識から証明方法を考え出すことができます。例えば、$(a + x)^{\frac{7}{3}}$ の展開は二項定理を使用する練習になります；$(a + x)^{\frac{7}{3}}$ に等しい級数を得た後、同じ方法で $(a + x)^{\frac{2}{3}}$ を展開します。これらの積は $\frac{7}{3} + \frac{2}{3} = 3$ より、$(a + x)^{3}$ の展開と同じであるべきであり、$a^{3} + 3a^{2}x + 3ax^{2} + x^{3}$ と等しいはずです。また、二項定理を使用して $(a + x)^{\frac{1}{2}}$ を展開した結果が $\sqrt{a + x}$ の平方根の計算結果と同じであることを試すこともできます。

さらに、どのような展開も得られた後は、その展開が元の式から導き出されたすべての性質を持っているかを確認すべきです。例えば、$\frac{1}{1 - x}$ は級数
\[
1 + x + x^{2} + x^{3} + \cdots, \quad \textit{無限に}
\]
と等しいことが証明されています。この式を $1 - x$ で掛けると $1$ になるべきであり、$1 - x^{2}$ で掛けると $1 + x$ になるべきです。すなわち
\[
\frac{1}{1 - x} \times (1 - x) = 1, \qquad
\frac{1}{1 - x} \times (1 - x^{2}) = 1 + x, \quad \textit{etc.}
\]
また、
\begin{align*}
a^{x} \&= 1 + x \Log a + \frac{x^{2} (\Log a)^{2}}{2} + \frac{x^{3} (\Log a)^{3}}{2 \cdot 3} + \cdots \textit{無限に}, \\
a^{y} \&= 1 + y \Log a + \frac{y^{2} (\Log a)^{2}}{2} + \cdots, \\
a^{x+y} \&= 1 + (x + y) \Log a + \frac{(x + y)^{2} (\Log a)^{2}}{2} + \cdots.
\end{align*}

今、$a^{x} × a^{y} = a^{x+y}$ から、最初の2つの級数の積は第3級数を与えるべきである。同様な他の多くの例が自然に思い浮かぶだろうし、それらに対する注意深い観察は、初学者にとってしばしば疑問視される一般性のある論理による各定理の証明を確認する助けとなるであろう。

\#\#\#\# 3. 特定の例と一般的なケース

III. 証明が符号の数と一般性によって困惑した場合、特定の例を選び、同じ証明法を適用させることで解決できる。例えば、二項定理が十分に明白でない場合でも、$(1 + x)^{\frac{2}{3}}$ または他の $(1 + x)^{\frac{m}{n}}$ のような特定の例に同様の論理を適用することができる。

また、積 $(x + a)(x + b)(x + c)$ 等、$n$ 個の因子を持つ一般形は、最初に2つの因子、次に3つ、そして4つを選んでから、一般的な積の形式を立証する論理を適用しようとする前に明らかにすることができる。

特に、順列と組合せの理論、および確率論がそのような関連性を持つように、これらも同様である。方程式の理論においては、最初に一般形
\[
x^{n} + Ax^{n-1} + Bx^{n-2} + \dots + Lx + M = 0,
\]
ではなく、3次または最多4次の方程式を選んで、そのような表現のすべての性質を立証することが望ましい。

しかし、これらの特定の例が扱われた後でも、一般的なケースは無視されず、全ての項を書き出すことができない不定数の項を持つ式に対する論理的推理の能力が鍛えられるべきである。これにより学生は、より高度な数学分野に進む際に成功への見通しを持てるようになるだろう。

\#\#\# まとめ

代数の学習では、各手続きを完全に理解し、その関連性を確認することが重要です。特定の例を通じて証明法を試すことで、一般的な理論がより深く理解できるようになります。また、多くの例題を解いて実践することで、代数的手続きの正確な適用が身につきます。これらの練習は時間と労力が必要ですが、その結果得られる深い洞察は、数学的な思考力を高め、新たな発見への扉を開くでしょう。

ご自身で手を動かし、計算を行い、疑問点を書き出して考えることで、代数の奥深さを体感してみてください。この旅を通じて、あなたが再び数学に挑戦する勇気と喜びを見つけることができるでしょう。

\end{document}
